\chapter{Le Fondamenta del Sistema Previdenziale: Teoria e Scelte Storiche}

Un sistema pensionistico non è mai neutro. Dietro ogni aliquota contributiva, ogni formula di calcolo, ogni soglia anagrafica, c'è una scelta su chi paga e chi riceve, su quale generazione sostiene il peso e quale lo scarica nel tempo. Per capire la crisi strutturale del sistema italiano occorre partire dalle fondamenta: le ragioni teoriche che giustificano l'intervento dello Stato, la logica matematica dei modelli di finanziamento, e le scelte storiche che hanno plasmato un'architettura previdenziale ancora gravata dalle sue contraddizioni originarie.

\section{Le ragioni dell'intervento pubblico}

\subsection{Funzione previdenziale, assicurativa e assistenziale}

Il sistema pensionistico pubblico svolge tre funzioni distinte che nella realtà si sovrappongono, ma che la teoria tiene accuratamente separate. La prima è la funzione previdenziale: garantire ai lavoratori in pensione un tenore di vita simile a quello raggiunto durante l'attività, a condizione che abbiano contribuito in misura adeguata. La seconda è la funzione assicurativa: proteggere i cittadini dal rischio di riduzione del reddito legato a eventi come l'invalidità, la morte prematura del capofamiglia, o semplicemente una longevità superiore alle aspettative. La terza è la funzione assistenziale: garantire un reddito minimo a chi non ha contribuito abbastanza per vivere in modo dignitoso, indipendentemente dalla storia lavorativa~\cite{Bosi2018, RosenGayerRapallini2023}.

Le tipologie di pensione corrispondono a queste funzioni in modo preciso~\cite{RosenGayerRapallini2023}:

\begin{center}
\begin{tabularx}{\textwidth}{|l|>{\raggedright\arraybackslash}X|>{\raggedright\arraybackslash}X|}
\hline
Tipologia & Beneficiario & Funzione prevalente \\
\hline
Vecchiaia & Ritirato per limiti di età & Previdenziale / assicurativa \\
\hline
Anzianità (anticipata) & Ritirato per anni di contribuzione & Previdenziale / assicurativa \\
\hline
Ai superstiti & Familiare del lavoratore deceduto & Assistenziale / previdenziale \\
\hline
Invalidità & Inabile al lavoro & Assicurativa / assistenziale \\
\hline
Sociale & Sprovvisto di reddito e incapace di lavorare & Assistenziale \\
\hline
\end{tabularx}
\end{center}

La distinzione rileva sul piano dell'analisi positiva perché ciascuno strumento ha una propria logica di finanziamento e di accesso. La pensione di invalidità risponde a un criterio medico, mentre quella di anzianità risponde a un criterio contributivo legato agli anni di versamenti accumulati dal lavoratore. Confondere i due piani produce distorsioni che si materializzano nei bilanci dell'ente previdenziale anni o decenni più tardi.

In Italia, tra gli anni Settanta e Ottanta, questi strumenti furono spesso impiegati al di fuori della loro logica originaria. Le pensioni di invalidità, in alcune aree del Mezzogiorno, vennero talvolta utilizzate come forma di sostegno al reddito in assenza di un vero reddito minimo garantito. Le pensioni di anzianità, rese accessibili da requisiti relativamente bassi introdotti con la riforma del 1969, divennero invece uno dei principali canali di prepensionamento industriale: nelle grandi imprese del Nord, le ristrutturazioni e le chiusure di reparti potevano essere accompagnate dall’uscita anticipata dei lavoratori, con una parte rilevante dei costi trasferita al sistema previdenziale~\cite{Bosi2018, Ferrera_1996}.

Ogni lavoratore mandato in pensione a cinquant'anni invece che a sessantacinque rappresentava quindici anni di prestazione aggiuntiva rispetto a quanto tecnicamente dovuto, senza una corrispondente copertura contributiva. Questo accumulo è quello che Bosi chiama ``debito previdenziale'': la differenza tra il valore attuale di tutte le pensioni promesse e il valore attuale di tutti i contributi futuri attesi. Una volta formatosi, è quasi impossibile ridurlo senza toccare i diritti acquisiti di chi già percepisce la pensione, oppure scaricando il costo sui lavoratori attivi attraverso aliquote più alte~\cite{Bosi2018}.

\subsection{Fallimenti del mercato e miopia individuale}

In presenza di mercati completi e di agenti pienamente razionali, l'intervento pubblico in materia previdenziale sarebbe difficile da giustificare sul piano dell'efficienza: per il primo teorema dell'economia del benessere, l'allocazione di mercato risulterebbe già Pareto-efficiente. L'individuo dell'ipotesi del ciclo vitale di Modigliani accumulerebbe risparmio durante la vita attiva per smobilizzarlo in vecchiaia, livellando il consumo nel tempo. La realtà si discosta da queste ipotesi su più fronti, e l'intervento pubblico risponde a fallimenti documentati~\cite{Bosi2018}.

I mercati assicurativi privati non riescono a offrire copertura completa contro il rischio di inflazione e contro i rischi sistemici, che per definizione non si diversificano, mentre la moltiplicazione di fondi privati in concorrenza genera costi di gestione e marketing che erodono il rendimento netto del pensionato~\cite{RosenGayerRapallini2023}. Vi è poi la miopia individuale, per cui gli esseri umani sistematicamente sottovalutano il futuro lontano e sovrastimano la propria capacità di risparmio volontario, al punto che senza un obbligo contributivo molti lavoratori arriverebbero alla vecchiaia senza riserve sufficienti~\cite{RosenGayerRapallini2023}.

L'azzardo morale completa il quadro. Un lavoratore che non risparmia sa che lo Stato non lo lascerà in miseria nella vecchiaia; questa convinzione genera un incentivo perverso che l'obbligatorietà del contributo risolve ex ante. A questo si aggiunge la selezione avversa: se la contribuzione fosse volontaria, i lavoratori più giovani e produttivi uscirebbero dal sistema lasciando solo i soggetti a maggior rischio, la base contributiva si restringerebbe e il sistema imploderebbe~\cite{RosenGayerRapallini2023}.

\section{I modelli di finanziamento a confronto}

\subsection{Il sistema a Ripartizione (PAYG): il patto intergenerazionale}

In un sistema a ripartizione i contributi versati oggi all'INPS escono immediatamente dall'ente per pagare le pensioni di chi è già in quiescenza. Il lavoratore di oggi finanzia il pensionato di oggi; quando raggiungerà a sua volta l'età pensionabile, sarà la generazione successiva a coprire il suo assegno. Il sistema regge finché ogni generazione accetta implicitamente questo accordo e finché le condizioni demografiche ed economiche non cambiano in modo tale da rendere impossibile l'adempimento~\cite{Bosi2018}.

Bosi definisce il meccanismo ``patto intergenerazionale'', perché il rendimento non dipende dai mercati finanziari ma dalla dimensione e dalla produttività della generazione successiva: una forza lavoro futura numerosa e produttiva sostiene pensioni adeguate, mentre una forza lavoro scarsa e a bassa crescita salariale comprime il monte contributivo disponibile~\cite{Bosi2018}.

Per formalizzare il meccanismo si usa il modello a generazioni sovrapposte, nella formulazione che da Samuelson e Diamond in poi costituisce il riferimento canonico per l'analisi dei sistemi a ripartizione~\cite{Samuelson1958, Diamond1965, Bosi2018}. Siano \(N_t\) gli anziani (pensionati) e \(N_{t+1}\) i giovani (lavoratori). La popolazione cresce al tasso costante \(n\), per cui \(N_{t+1} = (1+n)N_t\). Il parametro \(n\) misura il tasso di crescita della forza lavoro da un periodo al successivo: un valore \(n = 0{,}01\) indica che la nuova generazione di lavoratori è l'1\% più numerosa della precedente. Nelle economie mature contemporanee, con tassi di fertilità sotto il livello di rimpiazzo, \(n\) è prossimo a zero o negativo~\cite{Bosi2018, RosenGayerRapallini2023}.

Il salario per lavoratore è \(S_t\) e cresce al tasso di produttività \(m\). L'aliquota contributiva è \(c\). La pensione pro capite in ripartizione è~\cite{Bosi2018, RosenGayerRapallini2023}:

\[
P_{t+1}^{\text{ripartizione}} = c \cdot S_{t+1} \cdot \frac{N_{t+1}}{N_t} = c \cdot S_t (1+m)(1+n)
\]

Il tasso di rendimento implicito del sistema è pari a \(m+n\), ossia la somma del tasso di crescita della produttività e del tasso di crescita della popolazione occupata. In ripartizione, il rendimento non dipende dai mercati finanziari ma dall'andamento dell'economia reale~\cite{RosenGayerRapallini2023}.

Con \(c=0{,}33\), \(S=30.000\) euro, \(m=0{,}02\), \(n=0{,}01\) e un rapporto di 2 lavoratori per ogni pensionato, la pensione annua pro capite è:

\[
P = 0{,}33 \times 30.000 \times 1{,}02 \times 1{,}01 \times 2 = 20.430 \text{ euro}
\]

Se la demografia peggiora e \(n\) scende a \(-0{,}01\), il monte contributivo si contrae e ogni pensionato riceve di meno. Il sistema è strutturalmente esposto alla demografia: un calo della natalità, un allungamento della speranza di vita, una contrazione dell'occupazione riducono il monte contributivo o allungano il periodo di erogazione~\cite{RosenGayerRapallini2023}.

\subsection{Il sistema a Capitalizzazione (Fully Funded): rischi e opportunità}

Nel sistema a capitalizzazione ogni lavoratore accumula su un conto individuale i propri contributi, investiti sui mercati finanziari. Al pensionamento, la pensione equivale al montante accumulato rivalutato al tasso di rendimento \(r\), senza alcuna redistribuzione tra generazioni~\cite{Bosi2018, RosenGayerRapallini2023}:

\[
P_{t+1}^{\text{capitalizzazione}} = c \cdot S_t (1+r)
\]

Il confronto tra i due sistemi si riduce al confronto tra \(r\) e \((m+n)\). Il teorema di Samuelson-Aaron formalizza questa condizione~\cite{Samuelson1958, Aaron1966}: il PAYG è efficiente solo quando il tasso di crescita demografica ed economica supera stabilmente il rendimento reale del capitale. Nelle economie mature, con popolazioni stagnanti e rendimenti reali di lungo periodo sui mercati finanziari intorno al 4-5\%, la capitalizzazione sembrerebbe avvantaggiata. I rischi strutturali ne limitano però la superiorità~\cite{Bosi2018}.

Il rischio di mercato è il più evidente: crisi finanziarie come quella del 2008 possono azzerare decenni di accumulo senza meccanismi compensativi collettivi. Il rischio di longevità si aggiunge perché vivere più a lungo del previsto esaurisce il montante, e le rendite vitalizie offerte dal mercato sono costose proprio perché la longevità è un rischio aggregato non diversificabile~\cite{RosenGayerRapallini2023}.

Il rischio di inflazione merita un'analisi più articolata. L'effetto Fisher descrive il comportamento dei tassi nominali sui nuovi strumenti emessi dopo un aumento dell'inflazione. Il punto critico è che un fondo pensione è fatto in larga parte di titoli acquistati anni o decenni prima, quando l'inflazione era bassa e i tassi nominali erano corrispondentemente contenuti. Quelle posizioni sono aperte, i loro rendimenti sono fissi per contratto e non possono essere rivalutate retroattivamente: un portafoglio obbligazionario che rende il 3-4\% nominale subisce un'erosione reale significativa se l'inflazione sale al 7\% nei cinque anni precedenti il pensionamento. A questo si aggiunge un profilo istituzionale: nei mercati europei del dopoguerra, e in Italia durante la stagflazione degli anni Settanta, i tassi d'interesse erano soggetti a controlli politici che impedivano al meccanismo correttivo di attivarsi~\cite{RosenGayerRapallini2023}. I titoli indicizzati all'inflazione, come i BTP Italia o i CCT a cedola variabile, attenuano questo rischio per la quota di portafoglio che vi è investita, ma restano una componente minoritaria dei patrimoni previdenziali e non proteggono i titoli a cedola fissa già in portafoglio, sui quali la perdita reale resta integrale.

In ripartizione le pensioni sono spesso indicizzate all'inflazione per legge, e se i salari nominali crescono con l'inflazione crescono anche i contributi versati. Il rischio diventa collettivo e lo Stato agisce come garante residuale. Questo meccanismo cede tuttavia in caso di stagflazione: i contributi ristagnano perché i salari reali sono fermi, ma le pensioni indicizzate continuano a salire in termini nominali, producendo un disavanzo strutturale. L'effetto Fisher opera sui flussi futuri e non sugli stock passati, funziona imperfettamente nei mercati regolamentati, e non risolve il problema dell'incertezza individuale di lungo periodo~\cite{RosenGayerRapallini2023}.

\section{La genesi e l'evoluzione del sistema italiano}

\subsection{Dalle origini al dopoguerra: il passaggio alla ripartizione}

Il sistema pensionistico italiano nasce come sistema a capitalizzazione. Le prime pensioni di invalidità e vecchiaia furono istituite nel 1864 per i soli dipendenti statali; nel 1919, con il decreto legislativo 603, l'assicurazione obbligatoria fu estesa ai lavoratori dipendenti privati dell'industria, dell'agricoltura e dei servizi, gestita dalla Cassa Nazionale delle Assicurazioni Sociali~\cite{CERP2008}.

Due guerre mondiali e l'inflazione bellica e post-bellica distrussero le riserve accumulate. Con la legge 218 del 1952 si introduce il trattamento minimo di pensione e il sistema entra in una fase mista in cui il contributo pubblico diretto affianca i contributi dei lavoratori, operando già parzialmente con logica di ripartizione~\cite{CERP2008, Bosi2018}. Il passaggio definitivo avviene nel 1969, con la riforma Brodolini (legge 153/1969). La scelta non fu dettata soltanto da ragioni tecniche: i sindacati chiedevano pensioni più alte nell'immediato, e con la ripartizione bastava trasferire subito una quota più alta dei salari correnti, scaricando il costo sui lavoratori attivi anziché su riserve pregresse~\cite{Fondapol}.

L'Italia del 1969 aveva una struttura demografica favorevole. Come scrive Favero su Lavoce.info, nessuno fece i conti su cosa sarebbe successo quando quel rapporto si fosse invertito. La scelta di alzare subito le prestazioni correnti senza costituire coperture pluriennali aprì un disavanzo previdenziale crescente, nascosto dall'inflazione per un decennio e poi diventato visibile quando la crescita rallentò~\cite{LavoceRiformaDini}.

\subsection{Il metodo retributivo: logica, distorsioni e debito implicito}

L'abbandono della capitalizzazione coincise con l'adozione del metodo retributivo per il calcolo della pensione. La pensione non dipende dai contributi effettivamente versati, ma dalla retribuzione percepita negli ultimi anni di lavoro moltiplicata per un tasso legato agli anni di contribuzione~\cite{Bosi2018}:

\[
P = \tau \times W_{\text{pensionabile}}
\]

dove \(\tau\) è il tasso di rendimento (2\% per anno di contribuzione, da un minimo del 40\% per 20 anni a un massimo dell'80\% per 40 anni) e \(W_{\text{pensionabile}}\) è la retribuzione di riferimento rivalutata per inflazione più 1\%~\cite{RosenGayerRapallini2023}.

Il metodo retributivo scollega quasi completamente i benefici dai contributi versati, perché un lavoratore che ha guadagnato poco per trent'anni e molto negli ultimi cinque riceve una pensione ben superiore a quanto giustificato dai suoi versamenti complessivi, con un rendimento implicito sui contributi sistematicamente superiore al tasso di crescita dell'economia, in violazione del principio di sostenibilità finanziaria~\cite{LavoceRiformaDini}. Le prime avvisaglie di squilibrio emergono negli anni Ottanta: il rallentamento della crescita, il calo demografico e la riduzione del tasso di occupazione comprimono il monte contributivo, mentre le promesse retributive continuano a pesare sul bilancio INPS a ritmo crescente. Il debito previdenziale non compariva nei bilanci ufficiali, ma era reale come il debito esplicito, e la crisi degli anni Novanta avrebbe reso quel debito improvvisamente visibile~\cite{LavoceRiformaDini, Bosi2018}.

\subsection{La stagione delle riforme: dal metodo retributivo al sistema NDC}

La riforma Amato del 1992 (decreto legislativo 503/1992) è il primo intervento organico: aumenta l'età pensionabile, porta il requisito contributivo minimo per la pensione di vecchiaia da 15 a 20 anni, e amplia la finestra temporale di riferimento per il calcolo della retribuzione pensionabile. Si tratta di un intervento sul perimetro del sistema, senza toccare la logica di fondo del metodo retributivo~\cite{Bosi2018}.

La riforma strutturale arriva con il governo Dini nel 1995 (legge 335/1995), in un contesto di pressione per l'ingresso nell'euro con il debito pubblico già oltre il 120\% del PIL. L'obiettivo è sostituire il sistema retributivo con il sistema contributivo nozionale (NDC, Notional Defined Contribution): la pensione dipende dai contributi effettivamente versati durante l'intera vita lavorativa, capitalizzati a un tasso agganciato alla crescita del PIL. La scelta di usare la crescita del PIL come tasso di capitalizzazione ha una coerenza interna precisa: in un sistema a ripartizione, come dimostrato nel modello a generazioni sovrapposte, il tasso di rendimento implicito è proprio \((1+m)(1+n)\), ossia la crescita dell'economia. Usare quel tasso come coefficiente di accumulazione virtuale rendeva il sistema internamente coerente e, almeno in teoria, attuarialmente neutrale, nel senso che a parità di contributi versati il valore atteso della pensione è lo stesso per tutti e il sistema non trasferisce risorse né tra individui né tra generazioni oltre la crescita dell'economia~\cite{RosenGayerRapallini2023}.

La formula della pensione annuale lorda nel sistema NDC è:

\[
P = \alpha_L \cdot M_C^{NDC}
\]

dove \(\alpha_L\) è il coefficiente di trasformazione che dipende dalla speranza di vita residua all'età di pensionamento \(e_L\), e \(M_C^{NDC}\) è il montante contributivo, ottenuto come somma capitalizzata di tutti i contributi versati:

\[
M_C^{NDC} = \sum_{j=1}^{L} C_j \cdot (1+g)^{L-j}
\]

con \(g\) pari alla media mobile quinquennale del tasso di crescita del PIL nominale, \(C_j = c \cdot S_j\) e l'aliquota \(c = 0{,}33\) per i lavoratori dipendenti. Il coefficiente di trasformazione si ricava imponendo l'equilibrio finanziario tra montante e pensioni attese~\cite{RosenGayerRapallini2023}:

\[
M_C^{NDC} = \sum_{k=1}^{e_L} \frac{P}{(1+r_z)^{k-1}}
\quad \Longrightarrow \quad
P = \frac{M_C^{NDC}}{e_L}
\quad (\text{con } r_z = 0)
\]

Quanto più lunga è la speranza di vita al momento del pensionamento, tanto più piccolo è \(\alpha_L = 1/e_L\), e tanto più bassa è la pensione annua a parità di montante. Per il biennio 2023-2024 i coefficienti vigenti variano da un minimo del 4,270\% a 57 anni a un massimo del 6,655\% a 71 anni~\cite{RosenGayerRapallini2023}, e la revisione biennale li riduce ulteriormente al crescere della speranza di vita. È il meccanismo di aggiustamento automatico introdotto dalla riforma Dini e poi esteso dalla riforma Fornero ai requisiti di accesso~\cite{Bosi2018}; un'applicazione numerica completa del metodo, dal montante alla pensione netta, è già stata svolta nell'introduzione.

La riforma Dini applicava il sistema NDC integralmente solo a chi entrava nel mercato del lavoro dal 1 gennaio 1996. Chi aveva già 18 anni di contributi al 31 dicembre 1995 rimaneva interamente sotto il regime retributivo, anche per la parte maturata in seguito. Il risparmio effettivo sulla spesa pensionistica si sarebbe quindi materializzato solo decenni dopo, caricando sulle generazioni più giovani l'intero onere del riequilibrio e lasciando intatto il privilegio di chi era già profondamente radicato nel sistema~\cite{LavoceRiformaDini, Bosi2018}.

La risposta a questo gradualismo irrisolto arriva nel novembre 2011 con il decreto ``Salva Italia'' (decreto-legge 201/2011, convertito con legge 214/2011) del governo Monti, con lo spread a 570 punti base. La riforma Fornero estende il metodo contributivo a tutti i lavoratori a partire dal 1 gennaio 2012, unifica l'età pensionabile di vecchiaia a 67 anni per uomini e donne a partire dal 2019, e introduce il collegamento automatico dei requisiti di accesso alla speranza di vita: ogni due anni l'età richiesta per la pensione viene rivista in funzione delle variazioni della speranza di vita calcolate dall'Istat, così che a un allungamento della vita attesa corrisponda un posticipo dell'uscita dal lavoro. I risparmi stimati ammontavano a oltre 22 miliardi di euro entro il 2020~\cite{RepubblicaFornero2012}. La riforma si inserisce nel più ampio ciclo di interventi avviato nel 2004, al quale la Ragioneria Generale dello Stato attribuisce una riduzione del rapporto spesa/PIL di oltre 60 punti percentuali cumulati al 2060; alla sola riforma Fornero è ascrivibile circa un terzo di questo effetto~\cite{MEF_RGS2025_NdA, BosiGuerra2022}.

La risposta politica alla riforma si è tradotta in una serie di deroghe che hanno progressivamente eroso la coerenza attuariale del sistema. Quota 100 (decreto-legge 4/2019, triennio 2019-2021) consentiva il pensionamento con 62 anni di età e 38 di contributi. Quota 102, attiva per il solo 2022, aumentava i requisiti a 64 anni di età e 38 di contributi. Quota 103, introdotta per il 2023 e prorogata al 2024, ha fissato i requisiti a 62 anni di età e 41 di contribuzione, con la pensione calcolata integralmente con il metodo contributivo. Ciascuna di queste misure consente di percepire la pensione per più anni del previsto senza un corrispondente aumento del montante, scaricando il differenziale di costo sulla collettività e riproducendo, in forma attenuata, il meccanismo che la riforma Dini aveva cercato di interrompere~\cite{LavoceRiformaDini, MEF_RGS2025}.